%%%%%%%% MICE: Mixture of In-Context Experts for recurring concept drift. ICML template (icml2026.sty until the
%%%%%%%% 2027 style is released). Full draft of 2026-10-05. Every number is taken from logs/EXPERIMENT_LOG.md and
%%%%%%%% results_paper/ (code/make_paper_tables.py); tab_*_rows.tex are generated, do not edit them by hand.
%%%%%%%% The skeleton of 2026-10-03 is kept in old/main_skeleton_1003.tex.
\documentclass{article}

\usepackage{microtype}
\usepackage{graphicx}
\usepackage{booktabs}
\usepackage{hyperref}
\newcommand{\theHalgorithm}{\arabic{algorithm}}

\usepackage{icml2026}          % blind review
% \usepackage[accepted]{icml2026}

\usepackage{amsmath}
\usepackage{amssymb}
\usepackage{mathtools}
\usepackage{amsthm}
\usepackage[capitalize,noabbrev]{cleveref}
\usepackage{algorithm}
\usepackage{algorithmic}
\usepackage[textsize=tiny]{todonotes}

\theoremstyle{plain}
\newtheorem{theorem}{Theorem}
\newtheorem{proposition}{Proposition}
\newtheorem{lemma}{Lemma}
\theoremstyle{definition}
\newtheorem{definition}{Definition}
\newtheorem{assumption}{Assumption}
\newtheorem{observation}{Observation}
\theoremstyle{remark}
\newtheorem{remark}{Remark}

\input{math.tex}
\crefname{assumption}{Assumption}{Assumptions}
\crefname{observation}{Observation}{Observations}
\newcommand{\method}{\textsc{MICE}}   % Mixture of In-Context Experts
\makeatletter\newcommand{\rowinput}[1]{\@@input #1 }\makeatother   % expandable \input for generated table rows

\begin{document}

\twocolumn[
  \icmltitle{Recall or Relearn? A Mixture of In-Context Experts for Recurring Concept Drift}
  \icmlsetsymbol{equal}{*}
  \begin{icmlauthorlist}
    \icmlauthor{Anonymous Author(s)}{anon}
  \end{icmlauthorlist}
  \icmlaffiliation{anon}{Anonymous Institution}
  \icmlcorrespondingauthor{Anonymous Author(s)}{anon@anon.org}
  \icmlkeywords{concept drift, recurring concepts, tabular foundation models, in-context learning, online learning}
  \vskip 0.3in
]
\printAffiliationsAndNotice{}

\begin{abstract}
Tabular foundation models (TFMs) learn a data stream through their context, so a past concept can be recalled by
putting its rows back into the context, at no training cost. Current in-context stream learners do not use this:
a sliding context forgets a concept once it leaves the window, and a detector that resets the context is unsafe,
losing up to 8.4 points to a plain window on real streams. We ask when recalling a concept is worth more than
relearning it, and how to recall without paying when there is nothing to recall. We show that the value of memory is
set by the TFM's learning curve: the gain of recall over relearning is the area between the learning curve and the
accuracy of the stored context. We propose \method{}, a mixture of in-context experts, fast windows and stored
contexts, under a meta expert with two time scales. The fast level identifies a recurring concept within one batch
when the experts are separated by a margin; the slow level has a discounted regret of at most $\ln K/\eta$ against
the plain window at any time, with no assumption on the stream. On a pre-registered grid of recurring concepts, the
gain predicted from learning curves measured beforehand ranks the observed gains with Spearman 0.94, and \method{}
exceeds a reset detector on every stream with hard concepts, by 8.5 points on the hardest cell, and a sliding
window by 24 points on average. On 23 real streams \method{} is more accurate than all four in-context baselines on
18 and never more than 0.25 points below the plain window, including eight held-out segments evaluated once.
\end{abstract}

% =====================================================================================================
\section{Introduction}
\label{sec:intro}

A data stream revisits its past: seasons, weekdays, market regimes and sensor modes return. A learner that kept what
it knew about a concept can use it again when the concept recurs \citep{gama2014survey,lu2019learning}. For a
trained model this means keeping a pool of models and deciding which one to reuse
\citep{goncalves2013rcd,gama2014recurrent,zhao2025mooe}. Tabular foundation models (TFMs) such as TabPFN
\citep{hollmann2025tabpfn} and TabICL \citep{qu2025tabicl} change the cost of this idea. A TFM approximates posterior
prediction in one forward pass given a context of labelled rows \citep{muller2022pfn,nagler2023statistical}; its
knowledge of a concept \emph{is} a set of rows, and recalling the concept means conditioning on them. Memory needs no
training.

TFMs have already replaced incremental trees as stream learners: with a sliding memory they outperform adaptive tree
ensembles on non-stationary benchmarks \citep{lourenco2026streams}, and a prior with drifting causal models improves
them under temporal shift \citep{helli2024drift}. Neither line uses the free memory. A sliding context forgets a
concept once it leaves the window, and a drift-extrapolating prior models a trend rather than a return. Three
observations show what an in-context stream learner is missing, and what a remedy has to respect.

\begin{observation}[Sliding contexts forget]\label{obs:forget}
In the five batches after a concept recurs, a TFM with a FIFO context or with the dual memory of
\citet{lourenco2026streams} reaches 0.17, 0.39 and 0.51 accuracy on Sine, Stagger and Agrawal streams, whereas the
same TFM conditioned on the past rows of the recurring concept reaches at least 0.99; in steady state all contexts
are equally accurate.
\end{observation}

\begin{observation}[Relearning in context is cheap, for simple concepts]\label{obs:relearn}
A TFM relearns the concepts of the standard drift generators within about 200 rows, so resetting the context when a
drift detector fires \citep{gama2004ddm} already closes most of the gap of \cref{obs:forget} (0.81--0.88 in the same
five batches). Memory pays only for concepts that take many rows to learn: on mixtures with 100 centroids the TFM
needs more than 1000 rows to approach its asymptote.
\end{observation}

\begin{observation}[Resetting is unsafe]\label{obs:reset}
On 23 real streams the same reset detector is on average 1.2 points \emph{below} the plain FIFO context and up to 8.4
points below it: most rises of the error rate on real streams do not mean that the context has become useless.
\end{observation}

These observations raise the question of this paper: \emph{when should an in-context learner recall a past concept
rather than relearn it, and how can it do so without paying when there is nothing to recall?}

\textbf{Approach.} We treat every past context as an offline expert and the current window as an online expert, and
combine them with a meta expert, following the mixture of online and offline experts (MOOE) of
\citet{zhao2025mooe}. An in-context expert costs no training, which changes both the method and its analysis.
(i)~Which segments belong to the same concept is decided by a model-oriented discrepancy in the sense of
R-divergence \citep{zhao2023rdiv}; we show that the original form is blind to concepts that contradict each other on
the same inputs, which is the defining case of recurring drift, and that an excess-risk form repairs it.
(ii)~The meta expert has two time scales. A fast level restarts every batch and identifies a recurring concept in
one batch under a margin; a slow level decides whether following the fast level is worth it at all, with a
worst-case guarantee against the plain window. (iii)~The gain of recall over relearning then has a closed form in
the TFM's learning curve, so the value of memory can be measured before deployment.

\textbf{Contributions.}
\textbf{(1) Problem.} Recurring concept drift for in-context learners, where past contexts are free experts
(\cref{def:recur}), with three observations on forgetting, relearning and resetting.
\textbf{(2) Theory.} The blind spot of R-divergence under conflicting concepts and its excess-risk repair
(\cref{prop:blind}); identification in one batch under a margin (\cref{prop:block}); the recall-versus-relearn
formula (\cref{prop:recall}); and an any-time discounted regret bound for the slow level (\cref{prop:outer}).
\textbf{(3) Method.} \method{}, a training-free mixture of window and stored in-context experts whose every
component corresponds to one of these results.
\textbf{(4) Evidence.} Pre-registered tests on a controlled grid, where the learning curve predicts the gain of
memory and \method{} exceeds reset detection, window ensembles, incremental learners and trained MOOE; and 23 real
streams, where \method{} is the most accurate in-context policy on 18 and stays within 0.25 points of the plain
window on all.

% =====================================================================================================
\section{Recurring Concept Drift for In-Context Learners}
\label{sec:def}

\textbf{Setting.} A stream delivers batches of $B$ rows; the labels of a batch arrive one batch late. A frozen TFM
$f_\Theta(\cdot\mid\vx;\cC)$ predicts every batch from a context $\cC$ of at most $M$ labelled rows, and a
\emph{context policy} decides which rows form the context. The stream is a sequence of intervals $I_1,I_2,\dots$ with
distributions $P_1,P_2,\dots$ over $(\vx,y)$; concept drift occurs between intervals in the sense of
\citet{lu2019learning}.

\begin{definition}[Recurring in-context concept drift]\label{def:recur}
Let $\cC_k$ be a context of rows of interval $I_k$ and $f_\Theta(\cdot\,;\cC_k)$ the $k$-th \emph{in-context expert}.
The drift into interval $I_{K+1}$ is \emph{recurring} if $\min_{k\le K}\Div_T(P_{K+1}\,|\,P_k)\le\delta$, where
$\Div_T(Q\,|\,P)$ is the risk on $Q$ of the expert of $P$ minus that of the expert of $Q$ (\cref{def:excess}), and
\emph{new} otherwise.
\end{definition}

The definition is model-oriented: a drift is recurring when some stored context serves the new interval as well as
its own rows would. Three policies are the reference points of this paper. \emph{FIFO} keeps the last $M$ labelled
rows \citep{lourenco2026streams}. \emph{Reset} empties the context when the drift detection method (DDM) of
\citet{gama2004ddm} fires. \emph{Window ensemble} mixes FIFO contexts of several lengths. None of them has memory.

% =====================================================================================================
\section{Theory: From a Model-Oriented Discrepancy to Recall versus Relearn}
\label{sec:theory}

For a concept $P$ and a loss $\ell$, the risk of a context is
$\epsilon_P(\cC)=\E_{(\vx,y)\sim P}\,\ell\bigl(f_\Theta(\cdot\mid\vx;\cC),y\bigr)$, and the \emph{learning curve} of
$P$ is $R_P(n)=\E_{\cC\sim P^n}\,\epsilon_P(\cC)$. We write $p_\vx=P(\cdot\mid\vx)$, $P_X$ for the marginal and
$\mathrm{JS}$ for the Jensen--Shannon divergence. Proofs are in \cref{app:proofs}.

\subsection{Which Contexts Share a Concept}

R-divergence \citep{zhao2023rdiv} declares two distributions identical \emph{for a model} when the hypothesis learned
on their mixture has the same risk on both. For an in-context learner that hypothesis is the TFM conditioned on the
union of the two contexts:
\begin{equation}\label{eq:rdiv}
  \Div_R(P\|Q)=\bigl|\epsilon_P(\cC_P\cup\cC_Q)-\epsilon_Q(\cC_P\cup\cC_Q)\bigr| .
\end{equation}
For its population behaviour we use the idealisation behind prior-fitted networks
\citep{muller2022pfn,nagler2023statistical}: with contexts of equal size from $P$ and $Q$ the predictive
distribution tends to the conditional of the mixture $U=\tfrac12(P+Q)$,
$u_\vx=\alpha(\vx)p_\vx+(1-\alpha(\vx))q_\vx$ with $\alpha=P_X/(P_X+Q_X)$, and with a context from $P$ alone to
$p_\vx$. We denote these limits by $h_U$ and $h_P$.

\begin{definition}[Recall regret and excess-risk discrepancy]\label{def:excess}
\begin{align}
  \Div_T(Q\,|\,P) &= \epsilon_Q(h_P)-\epsilon_Q(h_Q), \label{eq:transfer}\\
  \Div_E(P,Q) &= \tfrac12\bigl[\epsilon_P(h_U)-\epsilon_P(h_P)\bigr] \notag\\
              &\quad+\tfrac12\bigl[\epsilon_Q(h_U)-\epsilon_Q(h_Q)\bigr].
  \label{eq:excess}
\end{align}
\end{definition}
$\Div_T$ is what recalling the stored context of $P$ costs while concept $Q$ is active; $\Div_E$ is what pooling the
two contexts costs each of them.

\begin{proposition}[Blind spot of R-divergence and its repair]\label{prop:blind}
Under the log-loss and the idealisation above:
\begin{enumerate}
  \item $\Div_E(P,Q)=\mathrm{JS}(P,Q)-\mathrm{JS}(P_X,Q_X)\ge0$, the part of the Jensen--Shannon divergence between
        the joints that covariate shift does not explain. If $P_X=Q_X$, then
        $\Div_E(P,Q)=\E_{\vx}\,\mathrm{JS}(p_\vx,q_\vx)$, which is zero if and only if $p_\vx=q_\vx$ almost surely.
  \item (Blind spot.) Let $P_X=Q_X$ and let both concepts be deterministic, $y=g_P(\vx)$ and $y=g_Q(\vx)$, with
        disagreement mass $\rho=\Pr_{\vx}[g_P(\vx)\neq g_Q(\vx)]$. Then $\Div_R(P\|Q)=0$ for every $\rho\in[0,1]$,
        whereas $\Div_E(P,Q)=\rho\ln2$. Under the 0-1 loss, $\Div_R=0$ again, while
        $\Div_T(Q\,|\,P)=\rho=2\,\Div_E(P,Q)$.
\end{enumerate}
\end{proposition}

R-divergence asks whether the mixture hypothesis treats the two sets \emph{equally}; two concepts that contradict
each other on the same inputs are treated equally badly, so the answer is yes although neither context helps the
other. The excess-risk form asks whether pooling \emph{costs} anything, and for the log-loss it is exactly the
concept-shift part of the divergence. With finite contexts a second term appears: if $P=Q$, pooling doubles the
context, and $\E[\epsilon_P(\cC_P\cup\cC_Q)-\epsilon_P(\cC_P)]=R_P(2n)-R_P(n)\le0$. The empirical excess is
interference minus pooling gain, so merging two contexts when it is at most $\delta$ merges exactly when pooling
helps more than the concepts interfere. \method{} uses the directed form \eqref{eq:transfer} as its merge statistic,
which is also the regret that recalling the stored expert would incur (\cref{def:recur}).

\subsection{The Fast Level: Identification under a Margin}

Following \citet{zhao2025mooe}, the regret of a mixture splits into the regret of the meta expert against the best
expert and the regret of that expert, $\Regret=\Regret_{\mathrm{ME}}+\Regret_{\mathrm{KE}}$. MOOE bounds
$\Regret_{\mathrm{ME}}\le\sqrt{T\ln K}$ for cumulative losses over an interval of $T$ rounds. After a recurrence the
cumulative loss still carries the previous concept, so the fast level of \method{} restarts every batch: weights
$w_k\propto\exp(-\eta\,\ell_k)$ from the log-loss $\ell_k$ of expert $k$ on the last labelled batch only. Its
guarantee comes from a margin.

\begin{lemma}[Margin]\label{lem:margin}
Let $w_k\propto\exp(-\eta\,\ell_k)$ over $K$ experts and $\ell_{k^\star}\le\ell_k-g$ for all $k\ne k^\star$. With
$\kappa=\ln\bigl(1+(K-1)e^{-\eta g}\bigr)$: (i) $w_{k^\star}\ge e^{-\kappa}$; (ii) on every row the log-loss of the
mixture is at most that of expert $k^\star$ plus $\kappa$; (iii) the mixture predicts the class of expert $k^\star$
on every row where its top two probabilities differ by more than $e^{\kappa}-1$.
\end{lemma}

\begin{proposition}[One batch of identification]\label{prop:block}
Consider a block of $J$ batches of one concept and log-losses clipped at $c$. If expert $k^\star$ leads every
labelled batch of the block by a margin $g$, then per row
$\Regret_{\mathrm{ME}}\le c+(J-1)(K-1)e^{-\eta g}$: one batch of identification, independent of the block length up
to an exponentially small term, against $\sqrt{JB\ln K}$ for the cumulative rule.
\end{proposition}

\subsection{When Recall Beats Relearning}

Let concepts recur in blocks of $L$ rows ($J=L/B$ batches) and let $a_k(n)$ be the accuracy learning curve of concept
$k$, non-decreasing in $n$. On the $c$-th visit of concept $k$ ($c\ge2$) its stored expert holds $n_c=\min(M,(c-1)L)$
rows. The \emph{relearning} learner resets its context at the block boundary; in MOOE's terms it is the online
expert, and the stored context is an offline expert that needed no training.

\begin{proposition}[Recall versus relearn]\label{prop:recall}
Assume the margin of \cref{prop:block} for the better of the stored expert and the window, with $e^\kappa-1$ below
the top-two gap of that expert. Then on the $c$-th visit the expected accuracy of the mixture minus that of
relearning, averaged over the block, is
\begin{equation}\label{eq:area}
  \Delta_{k,c}(L)=\frac1J\sum_{j=1}^{J-1}\Bigl[a_k(n_c)-a_k\bigl(\min(jB,M)\bigr)\Bigr]_+ ,
\end{equation}
the area by which the learning curve lies below the accuracy of the stored expert. It is zero when
$a_k(B)=a_k(M)$, at most $\tfrac{M}{L}\bigl(a_k(M)-a_k(B)\bigr)$, and zero on the first visit.
\end{proposition}

\Cref{eq:area} is $\Regret_{\mathrm{OE}}-\Regret_{\mathrm{KE}}$ in the decomposition of \citet{zhao2025mooe}: memory
helps when the area between the learning curve and the stored accuracy exceeds the cost of the meta expert, and
\cref{prop:block} makes that cost one batch, which the relearning learner loses as well. The value of memory is
therefore a property of the backbone and the concept that can be measured before deployment, and it explains
\cref{obs:relearn}: for concepts with a flat learning curve the area is zero.

\subsection{The Slow Level: A Safeguard at Any Time}
\label{sec:outer}

\Cref{lem:margin} speaks about streams on which one expert leads by a margin. Without a margin the batch-wise
restart follows noise: the expert that was best on the last batch need not be best on the next, which is the
situation of \cref{obs:reset}. \method{} therefore places a second meta expert above the first. It has two experts,
the \emph{default} (the FIFO window of $M$ rows, $k=0$) and the \emph{fast mixture} of \cref{lem:margin} ($k=1$),
and weights them on a slow time scale. Let $\ell_{t,k}\in[0,1]$ be the mean half Brier score
$\tfrac12\lVert p-e_y\rVert_2^2$ of expert $k$ on batch $t$, and
\begin{equation}\label{eq:outer}
\begin{aligned}
  L_{t,k}&=\gamma L_{t-1,k}+\ell_{t,k},\qquad v_{t,k}\propto\exp\bigl(-\eta_2\gamma L_{t-1,k}\bigr),\\
  p_t&=v_{t,0}\,p_{t,0}+v_{t,1}\,p_{t,1},
\end{aligned}
\end{equation}
with discount $\gamma\in(0,1]$ and $L_{0,k}=0$. For $\gamma=1$ this is the cumulative meta expert of MOOE with a
mixable loss; $\gamma<1$ forgets losses older than about $1/(1-\gamma)$ batches.

\begin{proposition}[Discounted regret of the slow level]\label{prop:outer}
Let $K$ experts be combined by \eqref{eq:outer} with $\eta_2\le\tfrac12$. For every sequence of labels and expert
predictions, every time $T$ and every expert $k$,
\[
  \sum_{t=1}^{T}\gamma^{\,T-t}\bigl(\ell_t(p_t)-\ell_{t,k}\bigr)\;\le\;\frac{\ln K}{\eta_2}.
\]
\end{proposition}

With $K=2$, $\eta_2=\tfrac12$ and $\gamma=0.9$ the bound is $2\ln2\approx1.39$ against a discounted horizon of
$1/(1-\gamma)=10$ batches: at every time, over the recent past, \method{} loses at most $0.14$ per batch in half
Brier score to the plain window \emph{and} to the fast mixture, whichever is better, with no assumption on the
stream. The discount makes this an any-time statement: after a change of regime the slow level needs about ten
batches, not the length of the stream, to follow the better of the two. The two levels divide the work. The fast
level identifies a recurring concept within one batch when there is something to identify (\cref{prop:block}); the
slow level decides whether identification is worth following. The bounded, exp-concave loss is essential: the step
is confined to $\eta_2\le\tfrac12$, and the unbounded log-loss lets one confidently wrong batch after a switch
dominate the slow level for the whole horizon (\cref{sec:ablation}).

% =====================================================================================================
\section{Method: A Mixture of In-Context Experts}
\label{sec:method}

\method{} (\cref{alg:mice}) runs a frozen TFM; no expert is trained. Each component implements one result of
\cref{sec:theory}.

\textbf{Experts.} \emph{Windows} are FIFO contexts of the last 100, 300 and $M$ labelled rows; they relearn fast and
cover concepts that are new or cheap to learn. \emph{Stored concepts} are contexts kept in a pool. Every 500 labelled
rows close a segment $S$. The segment joins the stored expert $\cC_k$ that minimises the estimated recall regret
$\widehat{\Div}_T(S\,|\,\cC_k)=\hat\epsilon_S(\cC_k)-\hat\epsilon_S(S)$ if this is at most $\delta=0.02$, where
$\hat\epsilon_S(S)$ is cross-fitted in two folds because an in-context learner memorises its context; otherwise it
founds a new stored expert (\cref{prop:blind}). A stored expert keeps its last $M$ rows, and the pool holds at most
12 experts (least recently used out).

\textbf{Fast level.} The windows and the stored experts are mixed with weights $w_k\propto\exp(-\eta\,\ell_k)$,
$\eta=10$, where $\ell_k$ is the log-loss on the most recent labelled batch (\cref{lem:margin}).

\textbf{Slow level.} The fast mixture and the $M$-row window are mixed by \eqref{eq:outer} with $\gamma=0.9$ and
$\eta_2=\tfrac12$ (\cref{prop:outer}).

\begin{algorithm}[t]
\caption{\method{} for one batch $t$}\label{alg:mice}
\begin{algorithmic}[1]
\STATE \textbf{Input:} unlabelled batch $X_t$; labels of batch $t-1$; windows $\mathcal W$; pool $\mathcal P$
\STATE update the windows with batch $t-1$; if 500 new labelled rows: close a segment $S$ and merge it into
       $\argmin_{k\in\mathcal P}\widehat{\Div}_T(S\,|\,\cC_k)$ if the minimum is at most $\delta$, else add $S$ to
       $\mathcal P$
\STATE $\ell_k\leftarrow$ log-loss of expert $k$ on batch $t-1$, for $k\in\mathcal W\cup\mathcal P$
\STATE $p^{\mathrm{fast}}\leftarrow\sum_k w_k\,f_\Theta(\cdot\mid X_t;\cC_k)$ with $w_k\propto e^{-\eta\ell_k}$
\STATE $L_k\leftarrow\gamma L_k+{}$half Brier score of $k\in\{\text{default},\text{fast}\}$ on batch $t-1$
\STATE \textbf{Output:} $v_0\,f_\Theta(\cdot\mid X_t;\cC_{\mathrm{default}})+v_1\,p^{\mathrm{fast}}$ with
       $v_k\propto e^{-\eta_2\gamma L_k}$
\end{algorithmic}
\end{algorithm}

The cost is one TFM call per expert and batch, at most 15, plus three calls when a segment closes. All constants
were set on synthetic development streams; none was tuned on a real stream.

% =====================================================================================================
\section{Experiments}
\label{sec:exp}

The experiments answer four questions. \textbf{RQ1}: does the learning curve predict the value of memory?
\textbf{RQ2}: does \method{} recall, that is, exceed relearning and window ensembles when concepts recur?
\textbf{RQ3}: is \method{} safe and accurate on real streams, where recurrence is weak or absent?
\textbf{RQ4}: what does each component contribute, and do the theoretical conditions hold?

\textbf{Protocol.} Prequential evaluation in batches of $B=100$ rows with labels one batch late; context budget
$M=1000$ for every policy; backbone TabPFN v2 with four ensemble members \citep{hollmann2025tabpfn}, and TabICL
\citep{qu2025tabicl} as a second backbone. \textbf{In-context baselines}: FIFO; the dual memory (LTM) of
\citet{lourenco2026streams}; DDM-triggered reset \citep{gama2004ddm}; a window ensemble of the three windows of
\method{} without the pool. \textbf{Other stream learners}: Adaptive Random Forest \citep{gomes2017arf}, Streaming
Random Patches \citep{gomes2019srp}, the Hoeffding Adaptive Tree \citep{bifet2009hat}, Leveraging Bagging
\citep{bifet2010levbag}, ADWIN Bagging, Learn++.NSE \citep{elwell2011nse}, and MOOE with trained experts
\citep{zhao2025mooe}, tuned on the development streams over 72 configurations and an extension of its grid
(\cref{app:details}). \textbf{Streams.} A controlled grid of recurring concepts: mixtures of 5, 30 or 100 Gaussian
centroids whose class assignment is permuted between three concepts, so that $P_X$ is fixed and concepts conflict;
blocks of 500 or 2000 rows; three cycles. Thirty streams from six standard generators (SEA, Stagger, Agrawal, Sine,
RBF, Hyperplane) with recurring concepts. Twenty-three real streams from 13 sources, the long ones in segments:
Electricity, seven INSECTS streams \citep{souza2020insects} (one in three segments), Forest CoverType, PokerHand and
Airlines (three segments each), Rialto, NOAA Weather, Spam and Phishing.

\textbf{Pre-registration.} Hypotheses, data and decision rules were written down before each test, and every test
stream was run once. The real streams were opened in four stages (\cref{tab:real}). The first three were also used,
after their one-time evaluation, to develop the slow level; the third held-out set, eight later segments of four
streams that no method had loaded, and the grid seeds 20 and 21 were evaluated after the two-level rule was frozen.
The pre-registered slow level used a step of 10; the step $\tfrac12$ required by \cref{prop:outer} was added to the
analysis before any result of \method{} on these sets was computed, and both are reported (\cref{app:prereg}).

\subsection{RQ1: The Learning Curve Predicts the Value of Memory}
\label{sec:rq1}

We measured the learning curves of the grid concepts, wrote down the gain predicted by \cref{prop:recall}, and only
then ran the streams. On seeds 10 and 11 the fast level alone exceeds reset on all eight streams with 30 or 100
centroids, more for shorter blocks and for harder concepts, as predicted, and the rank correlation between predicted
and observed gain over the six cells is 0.94; the magnitudes agree as well (0.55 versus 0.53 points, 3.0 versus 2.6,
10.4 versus 8.6). On the fresh seeds 20 and 21 with the full two-level method the correlation is 0.83
(\cref{tab:grid,fig:recall}). The six standard generators are the other end of the scale: their learning curves are flat after
about 200 rows, the predicted gain is at most 0.4 points, and the observed gain of \method{} over reset is 0.1 points
on average. Recurring-concept memory cannot be evaluated on the standard generators with a TFM; it needs concepts
that are costly to learn.

\begin{figure}[t]
  \centering
  \includegraphics[width=\columnwidth]{fig_recall.pdf}
  \caption{(a) Learning curves of the grid concepts with TabPFN v2, measured before the grid streams were run.
  (b) Gain of memory over reset per grid cell: predicted from these curves by \cref{prop:recall} against observed, for
  the fast level on seeds 10 and 11 and for \method{} on seeds 20 and 21; the dashed line is equality.}
  \label{fig:recall}
\end{figure}

\begin{table}[t]
  \centering
  \caption{Controlled grid, fresh seeds 20 and 21 (accuracy in \%, mean of two streams per cell). Win.: window ensemble; Fast: the
  fast level without the slow level. Pred.: gain over reset predicted by \cref{prop:recall} from learning curves
  measured beforehand; Gain: observed \method{} minus reset, in points.}
  \label{tab:grid}
  \footnotesize
  \setlength{\tabcolsep}{2.2pt}
  \begin{tabular}{rrccccccc}
    \toprule
    Cent. & Block & FIFO & Win. & Reset & Fast & \method{} & Pred. & Gain \\
    \midrule
\rowinput{tab_grid_rows.tex}
    \bottomrule
  \end{tabular}
\end{table}

\subsection{RQ2: Recall under Recurring Concepts}
\label{sec:rq2}

On the grid \method{} exceeds reset on all eight streams with 30 or 100 centroids and the window ensemble on all
eight (\cref{tab:grid}): the gain comes from memory, not from faster adaptation. Averaged over the twelve streams it
reaches 86.5 against 84.3 for reset, 81.8 for the window ensemble and 62.4 for FIFO; on the hardest cell the gain
over reset is 8.5 points. The slow level costs 0.4 points here against the fast level alone, the price of the
safeguard. Stream learners that train a model are far behind when a concept must be learned within a block: on seeds
10 and 11, ARF reaches 68.0, SRP 69.7 and the tuned MOOE 78.1, against 83.9 for reset and 86.2 for \method{}. On the
thirty generator streams \method{} matches reset (95.0 against 94.8), and the forgetting of \cref{obs:forget} is
repaired: 0.82--0.89 in the five batches after a recurrence against 0.17--0.51 for FIFO.

\subsection{RQ3: Real Streams}
\label{sec:rq3}

\begin{table*}[t]
  \centering
  \caption{Real streams: mean accuracy (\%) per evaluation stage, and the worst difference to FIFO over the streams
  of the stage (points). Every stage was run once when it was opened; the third held-out set was opened after the
  two-level rule was frozen. Per-stream results are in \cref{tab:realfull}.}
  \label{tab:real}
  \small
  \setlength{\tabcolsep}{4pt}
  \resizebox{\textwidth}{!}{%
  \begin{tabular}{lccccccc|cccc}
    \toprule
     & & \multicolumn{6}{c|}{Mean accuracy} & \multicolumn{4}{c}{Worst difference to FIFO} \\
    Stage & Streams & FIFO & LTM & Reset & Win.\ ens. & Fast only & \method{} & Reset & Win.\ ens. & Fast only & \method{} \\
    \midrule
\rowinput{tab_real_summary_rows.tex}
    \bottomrule
  \end{tabular}}
\end{table*}

Real streams rarely announce a recurrence, and the experts are close to each other. Here the requirement is
safety first. \Cref{tab:real} shows that every policy that reacts to the last errors pays for it somewhere: reset
loses up to 8.4 points to FIFO (PokerHand) and 1.2 on average, the window ensemble up to 1.4, and the fast level
alone up to 2.3. \method{} loses at most 0.25 points on any of the 23 streams, is above FIFO on 19 of them (+0.66 on
average), and is more accurate than all four in-context baselines on 18. Its largest gains are where the stream
revisits its past, Rialto (+4.7 points over FIFO, with a daily cycle) and the INSECTS streams (+0.2 to +1.9).

\textbf{Confirmatory set.} On the eight segments opened after the method was frozen, \method{} is within 0.08
points of FIFO on every segment, has the highest mean (85.2 against 84.9 for FIFO, 84.4 for the window ensemble and
82.3 for reset) and is the most accurate policy on five segments; the paired difference to FIFO over batches is
$+0.30$ points (block-bootstrap 95\% interval $[+0.20,+0.40]$). The pre-registered step of 10 gives 85.1, the same
worst case, and five of eight segments; the fast level alone is below FIFO by 1.7 and 0.7 points on the two
PokerHand segments, which the slow level turns into $+0.6$ and $+1.3$.

\textbf{Other learners and a second backbone.} On the 15 streams of the first three stages, the best trained
stream learner per stream (seven learners, MOOE in two configurations) is below \method{} on all 15, by at least 0.2
points, and below FIFO on 11. With
TabICL as the backbone and no constant changed, \method{} is the most accurate policy on all four streams of the
first held-out set (84.0 on average against 83.4 for FIFO and 83.5 for the window ensemble).

\subsection{RQ4: Components and Conditions}
\label{sec:ablation}

\textbf{Windows and pool.} Removing the pool, so that the two levels mix the three windows only, lowers accuracy on
every one of the 24 grid streams (by 5.6 points on average) and on 22 of the 23 real streams (by 0.35 points on
average, 2.1 on Rialto and 1.5 on a PokerHand segment): memory, not the choice of a window, carries the gain, on real
streams as well. Removing the two short windows instead lowers the generator streams from 95.0 to 93.6 and the grid
from 86.4 to 81.8 and leaves the real streams within 0.2 points on average; the pool forms exactly as many stored
experts as a generator stream has concepts. \textbf{Slow level.} It leaves the recall regime almost
unchanged (grid: 86.5 against 86.9 for the fast level alone) and removes the losses on real streams (worst case
$-0.25$ against $-2.3$; \cref{tab:real}). With the log-loss in place of the half Brier score, one grid stream falls
from 86.4 to 74.9: after a switch the fast mixture is confidently wrong for one batch, and an unbounded loss lets
that batch outweigh the next ten. \textbf{Discrepancy.} On 144 pairs of grid segments, the recall regret and the
excess-risk discrepancy separate ``different concept'' from ``same concept'' with AUROC 1.00, and the R-divergence
estimate with 0.85 (0.42 on the hardest cell), as \cref{prop:blind} predicts; on pairs of the same concept the
estimated excess is negative, down to $-0.2$ for 100 centroids, the pooling gain. \textbf{Margin.} On the grid the
leading expert holds weight 1.000 from the second batch of a recurring block, with a gap of 0.8--8.0 nats to the
runner-up, so $(K-1)e^{-\eta g}<0.005$ in \cref{prop:block}; on the real streams the median gap is 0.008--0.012
nats, the regime for which the slow level is built. \textbf{Regret.} Over 33 development streams the largest
discounted regret of the slow level is 0.15 against the window and 0.35 against the fast mixture, within the bound
of 1.39; on adversarial sequences the bound holds at $\eta_2=\tfrac12$ and fails for steps of 2 and 10
(\cref{app:details}).

% =====================================================================================================
\section{Related Work}
\label{sec:related}

\textbf{Concept drift and recurring concepts.} Drift detection and adaptation are surveyed by
\citet{gama2014survey} and \citet{lu2019learning}. Recurring concepts are handled by storing past models and
reusing them, selected by statistical tests \citep{goncalves2013rcd}, by meta-learners of each model's domain
\citep{gama2014recurrent} or by weighted ensembles of batch learners \citep{elwell2011nse}. MOOE
\citep{zhao2025mooe} combines offline experts of past intervals with an online expert under a meta expert with
regret guarantees; \method{} keeps its decomposition and replaces trained experts by stored contexts, the
cumulative meta expert by two time scales, and the generic regret by a quantity measured from the learning curve.
\textbf{Model-oriented discrepancy.} R-divergence \citep{zhao2023rdiv} and I-Div \citep{zhao2024idiv} measure
distribution discrepancy through a learned hypothesis; we identify the blind spot of the former under conflicting
concepts and use an excess-risk form to decide which contexts to pool. \textbf{TFMs on streams.}
\citet{lourenco2026streams} show that a TFM with a short- and long-term FIFO memory outperforms incremental
ensembles; Drift-Resilient TabPFN \citep{helli2024drift} retrains the prior to extrapolate temporal trends. Neither
models recurrence. Concept drift has likewise been redefined for other new model families, such as multi-modal
large language models \citep{yang2025adapting,yang2026turning}. \textbf{Online aggregation.} Exponential weights
with exp-concave losses and their logarithmic regret are classical \citep{cesabianchi2006prediction,hazan2007log};
\cref{prop:outer} is their discounted, any-time form for two in-context experts.

% =====================================================================================================
\section{Conclusion}
\label{sec:conclusion}

For an in-context learner, memory of a past concept is free to store and free to recall, and its value is the area
between the learning curve and the accuracy of the stored context. \method{} turns this into a method with one
component per result: an excess-risk discrepancy decides which contexts share a concept, a fast meta expert
identifies a recurring concept in one batch, and a slow meta expert guarantees, at any time and for any stream,
that following it costs little against the plain window. When concepts recur and are costly to learn, \method{}
exceeds reset detection, window ensembles and trained experts by a wide margin; on 23 real streams it is the most
accurate in-context policy on 18 and never more than 0.25 points below a sliding window, where reset detection
loses up to 8.4 points.

\textbf{Limitations.} The gains on real streams other than Rialto are below two points; the claim there is safety
with a small average gain, and the large gains are for recurring concepts with steep learning curves.
The confirmatory real segments come from four streams whose earlier segments were used in development, and every
stream has one run. \Cref{prop:outer} bounds the half Brier score rather than the accuracy, and
\cref{prop:blind,prop:recall} hold for the idealised predictive and under a margin. \method{} calls the TFM once per
expert, 5--18 times the time of a single context in our runs.

\bibliography{refs}
\bibliographystyle{icml2026}

% =====================================================================================================
\newpage
\appendix
\onecolumn

\section{Proofs}
\label{app:proofs}

\begin{proof}[Proof of \cref{prop:blind}]
For the log-loss, $\epsilon_P(h)=\E_{P_X}\bigl[H(p_\vx)+\KL(p_\vx\|h_\vx)\bigr]$ for any conditional $h$, with $H$ the
entropy. Hence $\epsilon_P(h_P)=\E_{P_X}H(p_\vx)$ and $\epsilon_P(h_U)-\epsilon_P(h_P)=\E_{P_X}\KL(p_\vx\|u_\vx)$, and
likewise for $Q$. By the chain rule, $\KL(P\|U)=\KL(P_X\|U_X)+\E_{P_X}\KL(p_\vx\|u_\vx)$ and the same for $Q$;
averaging the two identities gives $\mathrm{JS}(P,Q)=\mathrm{JS}(P_X,Q_X)+\Div_E(P,Q)$, and $\Div_E\ge0$ as a mean of
KL terms. If $P_X=Q_X$ then $\alpha\equiv\tfrac12$, $u_\vx=\tfrac12(p_\vx+q_\vx)$ and
$\Div_E=\E_\vx\mathrm{JS}(p_\vx,q_\vx)$, which vanishes if and only if $p_\vx=q_\vx$ almost surely. This proves (1).
For (2), $\Div_R(P\|Q)=\bigl|\E_{P_X}[H(p_\vx)+\KL(p_\vx\|u_\vx)]-\E_{Q_X}[H(q_\vx)+\KL(q_\vx\|u_\vx)]\bigr|$.
Deterministic concepts have $H(p_\vx)=H(q_\vx)=0$, and $u_\vx$ puts mass $\tfrac12$ on each of $g_P(\vx)$ and
$g_Q(\vx)$ where they differ and mass $1$ where they agree, so
$\KL(p_\vx\|u_\vx)=\KL(q_\vx\|u_\vx)=\ln2\cdot\mathbf 1[g_P(\vx)\neq g_Q(\vx)]$; the two terms cancel and
$\Div_E=\rho\ln2$. Under the 0-1 loss, with ties of $u_\vx$ broken independently of the source of the query, the
prediction at a disagreeing $\vx$ is $g_P(\vx)$ or $g_Q(\vx)$ with probability $\tfrac12$ each, so each risk is
$\rho/2$ and $\Div_R=0$; $h_P$ errs on $Q$ exactly where the concepts disagree, so $\Div_T(Q\,|\,P)=\rho$.
\end{proof}

\begin{proof}[Proof of \cref{lem:margin}]
(i) $w_{k^\star}=\bigl(1+\sum_{k\ne k^\star}e^{-\eta(\ell_k-\ell_{k^\star})}\bigr)^{-1}\ge(1+(K-1)e^{-\eta g})^{-1}$.
(ii) $-\ln\sum_kw_kp_k(y)\le-\ln\bigl(w_{k^\star}p_{k^\star}(y)\bigr)\le-\ln p_{k^\star}(y)+\kappa$.
(iii) For classes $y_1,y_2$,
$\sum_kw_k\bigl(p_k(y_1)-p_k(y_2)\bigr)\ge w_{k^\star}\bigl(p_{k^\star}(y_1)-p_{k^\star}(y_2)\bigr)-(1-w_{k^\star})$,
which is positive when the gap of $k^\star$ exceeds $(1-w_{k^\star})/w_{k^\star}\le e^{\kappa}-1$.
\end{proof}

\begin{proof}[Proof of \cref{prop:block}]
The first batch of the block is predicted with weights from a batch of the previous concept; its per-row regret is at
most $c$. Every later batch is predicted with weights from a labelled batch of the block, on which $k^\star$ leads by
$g$, so \cref{lem:margin}(ii) bounds its per-row regret by $\ln(1+(K-1)e^{-\eta g})\le(K-1)e^{-\eta g}$.
\end{proof}

\begin{proof}[Proof of \cref{prop:recall}]
Both learners lose the first batch, whose labels have not arrived. Before batch $j+1$ the relearning context holds
$\min(jB,M)$ rows of the block, so its expected accuracy is $a_k(\min(jB,M))$; the stored expert has expected
accuracy $a_k(n_c)$. By \cref{lem:margin}(iii) the mixture predicts as the better of the two, which gives the
positive part. Each term is at most $a_k(M)-a_k(B)$ and vanishes for $jB\ge M$, so at most $M/B$ of the $J$ terms are
non-zero. On the first visit there is no stored expert of the concept.
\end{proof}

\begin{proof}[Proof of \cref{prop:outer}]
The half Brier score $\ell(p)=\tfrac12\lVert p-e_y\rVert^2$ has gradient $p-e_y$ and Hessian $I$, and
$\lVert p-e_y\rVert^2\le2$ on the simplex, so $\nabla^2\ell\succeq\eta_2\,\nabla\ell\nabla\ell^{\top}$ for
$\eta_2\le\tfrac12$: $\ell$ is $\eta_2$-exp-concave. The mean over the rows of a batch is $\eta_2$-exp-concave as
well, since its Hessian is the mean Hessian and $\tfrac1n\sum_ig_ig_i^\top\succeq\bar g\bar g^\top$ for the row
gradients $g_i$ with mean $\bar g$. Hence $e^{-\eta_2\ell_t(p_t)}\ge\sum_kv_{t,k}e^{-\eta_2\ell_{t,k}}$. Let
$W_t=\sum_ke^{-\eta_2L_{t,k}}$, $W_0=K$. Then
$W_t=\bigl(\sum_ke^{-\eta_2\gamma L_{t-1,k}}\bigr)\sum_kv_{t,k}e^{-\eta_2\ell_{t,k}}$, and by the power-mean inequality
$\sum_kx_k^{\gamma}\le K^{1-\gamma}\bigl(\sum_kx_k\bigr)^{\gamma}$ for $x_k\ge0$,
\[
  \ln W_t\;\le\;(1-\gamma)\ln K+\gamma\ln W_{t-1}-\eta_2\,\ell_t(p_t).
\]
Unrolling gives $\ln W_T\le\ln K-\eta_2\sum_{t\le T}\gamma^{T-t}\ell_t(p_t)$, because
$(1-\gamma)\sum_{t\le T}\gamma^{T-t}+\gamma^{T}=1$. On the other hand
$\ln W_T\ge-\eta_2L_{T,k}=-\eta_2\sum_{t\le T}\gamma^{T-t}\ell_{t,k}$.
\end{proof}

\section{Pre-Registration Record}
\label{app:prereg}

All hypotheses below were written to the experiment log before the corresponding streams were run, and every test
stream was run once; interrupted runs were resumed from their checkpoints without rerunning a finished cell.

\textbf{Grid, seeds 10 and 11} (fast level alone). Hypotheses: gain over reset positive on all eight streams with 30
or 100 centroids, larger for shorter blocks and for more centroids; Spearman correlation between predicted and
observed gain at least 0.8; the window ensemble below the method on these eight streams. All held (correlation
0.94). An incidental expectation of no gain with 5 centroids was missed on one stream, where the reset baseline
detected late.

\textbf{Generator streams, seeds 10--14, and five real development streams.} Non-inferiority to reset on the
generator streams ($+0.50$ points, 25 of 30 positive) and on each real stream held, as did the replication of
\cref{obs:forget}. An exploratory expectation that the INSECTS stream with re-occurring drift would gain most was not
met (it ranked second).

\textbf{First held-out set} (four streams; fast level alone). Non-inferiority to reset on every stream and highest
mean accuracy held; the method was the most accurate on three of four streams. A replication with TabICL held as
well.

\textbf{Second held-out set} (six streams; fast level alone). Non-inferiority to reset held. The hypothesis that the
fast level alone would be the most accurate method on at least four of six streams was \emph{not} met (one of six),
and on PokerHand it was 2.3 points below FIFO. This result motivated the slow level; all real streams seen so far
count as development data from here on.

\textbf{Third held-out set} (eight segments; two-level rule frozen). Two candidates for the slow level were
registered before any result existed: the log-loss with step 1, and the half Brier score with step 10, added after
the log-loss failed on a cached synthetic test stream. For both, non-inferiority to FIFO and reset (tolerance 0.5
points) and to the fast level alone held on every segment, and the mean was the highest. ``Most accurate on at least
five of eight segments'' held for the half Brier score (five) and not for the log-loss (four). The step
$\eta_2=\tfrac12$ with tempered weights, the method of this paper, was added to the analysis after the baselines of
three segments were visible and before any result of a two-level rule was computed; it was not part of the
registered hypotheses. It gives five of eight segments and the highest mean.

\textbf{Grid, seeds 20 and 21} (half Brier score, step 10). Mean difference to the fast level alone at least $-0.5$
points ($-0.14$), above reset on all eight streams with 30 or 100 centroids, and within one point of reset with 5
centroids: all held. The step $\tfrac12$ is 0.23 points lower on average and satisfies the same three statements.

\section{Experimental Details and Full Results}
\label{app:details}

\textbf{Streams.} The grid and generator streams are produced by \texttt{code/grid\_streams.py} and
\texttt{code/streams.py}; generator streams have blocks of 2000 rows and three cycles over three or four concepts.
Real streams with more than 100{,}000 rows are cut into segments of 100{,}000 (CoverType and one INSECTS stream:
150{,}000 for the first segment). The second held-out set consists of all streams of a public collection of drift
data sets that have at least 5000 rows, at most ten classes and had not been used before.
\textbf{Baselines.} LTM follows Algorithm~1 of \citet{lourenco2026streams} with a short-term share of 0.75. MOOE
follows the original algorithm with linear experts on random Fourier features, generalised to several classes; it
was tuned on the five development streams over feature type, interval length, step size and regularisation (72
configurations), and its grid was extended twice because the optimum lay on the boundary; the best configuration of
every round is evaluated. The incremental learners use the default settings of \texttt{river} \citep{montiel2021river}
with ten base models and the same batch protocol.
\textbf{Regret under adversarial sequences.} With synthetic expert predictions and labels chosen to hurt the mixture
or the current leader ($K\in\{2,5\}$, $\gamma\in\{0.5,0.9,0.99,1\}$, 400 rounds), the discounted regret at
$\eta_2=\tfrac12$ stays below $\ln K/\eta_2$ in every configuration (at most 0.67 for $K=2$), exceeds it at
$\eta_2=2$ for $\gamma\ge0.9$, and reaches 4.4 at $\eta_2=10$, $\gamma=0.9$.

\begin{table}[h]
  \centering
  \caption{Real streams: accuracy (\%) per stream. Fast only: the fast level without the slow level; $\eta_2=10$: the
  pre-registered step; Trained: the best trained stream learner (not run on the third held-out set);
  $\Delta$: \method{} minus FIFO in points.}
  \label{tab:realfull}
  \small
  \begin{tabular}{lccccccccc}
    \toprule
    Stream & FIFO & LTM & Reset & Win.\ ens. & Fast only & $\eta_2=10$ & \method{} & Trained & $\Delta$ \\
    \midrule
\rowinput{tab_real_rows.tex}
    \bottomrule
  \end{tabular}
\end{table}

\end{document}
